\textbf{结论名称：} 切线长定理（过圆外一点作圆的两条切线，切线段长相等）

\textbf{适用条件：} 点 \(P\) 在 \(\odot O\) 外，\(PA\)、\(PB\) 分别切圆于 \(A\)、\(B\)。

\textbf{变量说明：} 设圆心为 \(O\)，半径为 \(r\)，\(OP = d\;(d>r)\)，切线长记为 \(l = PA = PB\)。

\textbf{核心公式：}
\[
\boxed{PA = PB = \sqrt{OP^{2} - r^{2}} = \sqrt{d^{2} - r^{2}}}
\]
且 \(OP\) 平分 \(\angle APB\)，\(OP\) 垂直平分切点弦 \(AB\)。

\textbf{使用场景：} 求切线长；证明线段相等；处理切点弦及与圆幂相关的长度计算；解析几何中列切线方程并求切线长。

\textbf{简单例子：} 已知 \(\odot O\) 半径 \(r=3\)，点 \(P\) 满足 \(OP=5\)，则切线长
\[
PA = PB = \sqrt{5^{2} - 3^{2}} = 4.
\]
利用该结论可直接由勾股定理得到，无需建立坐标系。